\documentclass[11pt, oneside]{article} 
\usepackage{geometry}
\geometry{letterpaper} 
\usepackage{graphicx}
	
\usepackage{amssymb}
\usepackage{amsmath}
\usepackage{parskip}
\usepackage{color}
\usepackage{hyperref}

\graphicspath{{/Users/telliott/Dropbox/Github-Math/figures/}}
% \begin{center} \includegraphics [scale=0.4] {gauss3.png} \end{center}

\title{Steiner-Lehmus revisited}
\date{}

\begin{document}
\maketitle
\Large

%[my-super-duper-separator]

We go back to the Steiner Lehmus \hyperref[sec:SLtheorem]{\textbf{theorem}} to look at some proofs that depend on trigonometry including one that is highly algebraic and invokes Stewart's theorem.

\subsection*{simple trigonometric proof}

This proof is limited to angles less than a right angle, but it has the advantage of being particularly simple.  (Of course, in an isosceles triangle, the base angles must be acute, but we are going to assume that one is larger than the other).

\begin{center} \includegraphics [scale=0.35] {lehmus_base.png} \end{center}

We are given that $m=n$ and the angles are bisected.  Calculate (twice) the total area of the triangle in two different ways:
\[ 2 \mathcal{A} = mc \sin \alpha + mb \sin \alpha = nc \sin \beta + na \sin \beta \]

Aiming for a contradiction, let $\alpha > \beta$.  

Then we can deduce that 
\[ \sin \alpha > \sin \beta \]
Since, for $\theta$ in $[0,\pi/2]$, the greater angle has the larger sine.  And since $m = n$
\[ mc \sin \alpha > nc \sin \beta \]
and so referring back to the equality, it must be that
\[ mb \sin \alpha < na \sin \beta \]

From the law of sines:
\[ \frac{a}{\sin A} = \frac{b}{\sin B} = \frac{c}{\sin C} \]
\[ a = c \ \frac{\sin A}{\sin C} \]
\[ b = c \ \frac{\sin B}{\sin C} \]
and then substituting:
\[ mc \ \frac{\sin B}{\sin C} \sin \alpha < nc \ \frac{\sin A}{\sin C} \sin \beta \]

But $m = n$ so cancel to obtain
\[ \sin B \sin \alpha < \sin A \sin \beta \]

Employ the double angle formula and cancel the factor of $2$:
\[ \sin \beta \cos \beta \sin \alpha < \sin \alpha \cos \alpha \sin \beta \]
\[ \cos \beta < \cos \alpha \]

Recall that we have $\alpha > \beta$.  But for $\theta$ in $[0,\pi/2]$, the greater angle has the smaller cosine.

This is a contradiction.  Letting $\alpha < \beta$ leads to a contradiction as well.

Therefore $\alpha = \beta$.

[adapted from \url{https://dc.etsu.edu/etd/1169/} and referenced there as Seydel and Newman (1983) which is, unfortunately, paywalled.  ]

\subsection*{another look at Coxeter's lemma}

The proof in a previous chapter from Coxeter used what he calls Lemma 2.  

Here is a proof of the same statement based on the law of sines (\url{https://arxiv.org/pdf/2601.15591}).  It is closely related to the one just given above.

\emph{Lemma}.

If a triangle has two angles of different measure, the larger angle has the smaller internal bisector.
\begin{center} \includegraphics [scale=0.35] {SL_lemma.png} \end{center}

Let $\triangle ABC$ have bisected base angles $A$ and $B$ with bisector lengths $l$ and $l'$.  Suppose $A > B$.  Then $l < l'$.

Let the exterior angle $y = \gamma + \alpha$.

\emph{Proof}.

Reflect $l'$ so it originates from vertex $A$ (right panel).

In the blue triangle the angle at $A$ is $\beta$.  Let the difference $\alpha - \beta = \delta$.

In the same triangle, the total $\angle ABE$ is $2 \alpha = 2 \beta + 2 \delta$.

The angle where $l$ meets $BC$ is $y$ as before.

\begin{center} \includegraphics [scale=0.35] {SL_lemma.png} \end{center}

We claim that $x=y$ because both are supplementary to angles of the same measure, namely, $3 \beta + \delta = 2 \beta + \alpha$.

By the Law of Sines
\[ \frac{\sin x}{c} = \frac{\sin 2 \beta + \delta}{AD} \]

But also by the same law
\[ \frac{\sin y}{c} = \frac{\sin 2 \beta}{l} \]

equating the two (since $x=y$) we have
\[ AD = \frac{\sin 2 \beta + \delta}{\sin 2 \beta} \cdot  l \]

The larger angle has the larger sine, so the factor is larger than 1 which means that $AD > l$.  

But $AE = l' > AD$ so $l' > l$.

The larger angle has the smaller bisector.

$\square$


\subsection*{6.  Trigonometric proof by Paul Yiu}

Let the vertices and sides be labeled in the usual way.  Let the half-angles be $\beta$ and $\gamma$ also.

Let $AC$ (side $b$) be divided by $BY$ into $u$ and $U$, and let $AB$ (side $c$) be similarly divided by $CZ$ into $v$ and $V$.

\begin{center} \includegraphics [scale=0.18] {Steiner_Lehmus_Yiu.png} \end{center}

We will show that assuming $\beta < \gamma$ will lead to a contradiction:
\[ \frac{b}{u} < \frac{c}{v} \ \ \ \ \text{and also} \ \ \ \  \frac{b}{u} > \frac{c}{v} \]

First consider

\[ \frac{b}{u} - \frac{c}{v} = \frac{u + U}{u} - \frac{v + V}{v} \]
\[ = \frac{U}{u} - \frac{V}{v} \]
By the angle bisector theorem $a/U = c/u$ so $U = au/c$ and similarly $V = av/b$.  So
\[ = \frac{a}{c} - \frac{a}{b} \]
Since $c > b$ (by Euclid I.19, since $\gamma > \beta$) this expression is less than zero.  Retrieving the original LHS:
\[ \frac{b}{u} - \frac{c}{v} < 0 \]
\[ \frac{b}{u} < \frac{c}{v} \]

The second part is

\[ \frac{b}{u} \div \frac{c}{v} = \frac{b}{c} \ \frac{v}{u} \]
By the Law of Sines
\[ = \frac{\sin B}{\sin C} \ \frac{v}{u} \]
\begin{center} \includegraphics [scale=0.18] {Steiner_Lehmus_Yiu.png} \end{center}
By the sum of angles (really double angle)
\[ = \frac{2 \cos \beta \sin \beta}{2 \cos \gamma \sin \gamma} \ \frac{v}{u} \]
\[ = \frac{\cos \beta}{\cos \gamma} \ \frac{\sin \beta}{u} \ \frac{v}{\sin \gamma} \]
Drop the altitude $h$ from $Y$ to $AB$ (side $c$).  Then
\[ \frac{h}{BY} = \sin \beta \]
\[ \frac{h}{u} = \sin A \ \ \ \ \ \ \ \ \frac{h}{\sin A} = u  \]
Thus $\sin \beta/u = \sin A/BY$:
\[ = \frac{\cos \beta}{\cos \gamma} \ \frac{\sin A}{BY} \ \frac{CZ}{\sin A} \]
The third term has been massaged in just the same way.

And then, almost everything cancels!
\[ = \frac{\cos \beta}{\cos \gamma} > 1 \]
The inequality comes from the assumption that $\beta < \gamma$.  Retrieve the original LHS:
\[ \frac{b}{u} \div \frac{c}{v} > 1 \]
\[ \frac{b}{u} > \frac{c}{v} \]

We have that
\[ \frac{b}{u} < \frac{c}{v} \ \ \ \ \text{and also} \ \ \ \  \frac{b}{u} > \frac{c}{v} \]

This is a contradiction.  It cannot be that $\beta < \gamma$.  Similar logic reaches a contradiction for $\beta > \gamma$, so it must be that $\beta = \gamma$.  Then it is trivial to show that $\triangle ABC$ is isosceles.

$\square$

\subsection*{7.  Algebraic proof of Steiner-Lehmus}

Here is a last proof of the theorem.  First we revisit Stewart's theorem in the case of the bisected angle.  Substitute the symbol $t$ (for transversal).

\begin{center} \includegraphics [scale=0.15] {steiner1.png} \end{center}
\[ ab = t^2 + xy \]
\[ t^2 = ab - xy \]

This is a general formula for the length of the transversal when it is an angle bisector.  We need to eliminate the $xy$ term.  One way to write the fundamental relationship for a bisected angle is
\[  \frac{a}{b} = \frac{y}{x} \]
\[  \frac{a+b}{b} = \frac{x+y}{x} \]
\[ x = \frac{b(x+y)}{a+b} = \frac{bc}{a + b} \]

By symmetry 
\[ y = \frac{a}{b} x = \frac{ac}{a + b} \]
So
\[ xy = \frac{abc^2}{(a+b)^2} \]

And then
\[ t^2 = ab - \frac{abc^2}{(a+b)^2} \]
The length (squared) of the transversal to side $c$ is symmetric in $a$ and $b$, as it should be.

\[ = ab \ [ \ 1 - \frac{c^2}{(a+b)^2} \ ] \]
\[ = ab \ [ \ \frac{(a + b)^2 - c^2}{(a+b)^2} \ ] \]
\[ = ab \ [ \ \frac{(a + b + c)(a + b - c)}{(a+b)^2} \ ]  \]

We can apply this formula to the main problem, since there are two transversals which are equal.  The algebra gets pretty messy.

As an overview, the key step is to produce (something like):
\[ (b-a) \ [ \ 3 abc + ab(a + b) + c^3 + c^2(a+b) \ ] \ = 0 \]

Since $a,b,c$ are all positive, we have $(b-a)$ times something which is positive.  Thus $b - a = 0$ and so $a = b$.

Rewrite the $t^2$ expression for sides $a$ and $b$:
\[ t_a^2 = bc \ [ \ \frac{(a + b + c)(-a + b + c)}{(b+c)^2} \ ]  \]
\[ t_b^2 = ac \ [ \ \frac{(a + b + c)(a - b + c)}{(a+c)^2} \ ]  \]

Set them equal and cancel some terms:
\[ b \ [ \ \frac{(-a + b + c)}{(b+c)^2} \ ] = a \ [ \ \frac{(a - b + c)}{(a+c)^2} \ ]  \]
\[ b(b + c - a)(a + c)^2 = a(a + c - b)(b+c)^2 \]

We will obtain nine terms on each side.  Start with the LHS
\[ b(b + c - a)(a + c)^2 \]
\[ = (b^2 + bc - ab)(a^2 + 2ac + c^2) \]
\[ = a^2b^2 + 2ab^2c + b^2c^2 + a^2bc + 2abc^2 + bc^3 - a^3b - 2a^2bc - abc^2 \]
\[ = a^2b^2 + 2ab^2c + b^2c^2 + abc^2 + bc^3 - a^3b - a^2bc \]

Now for the RHS (we \emph{could} do a symbol swap with $b$ for $a$ and $a$ for $b$, but ...):
\[ a(a + c - b)(b+c)^2 \]
\[ = ( a^2 + ac - ab)(b^2 + 2bc + c^2) \]
\[ = a^2b^2 + 2a^2bc + a^2c^2 + ab^2c + 2abc^2 + ac^3 - ab^3 - 2ab^2c - abc^2 \]
\[ = a^2b^2 + 2a^2bc + a^2c^2 + abc^2 + ac^3 - ab^3 - ab^2c \]

Compare LHS with RHS for symmetry
\[ = a^2b^2 + 2ab^2c + b^2c^2 + abc^2 + bc^3 - a^3b - a^2bc \]
\[ = a^2b^2 + 2a^2bc + a^2c^2 + abc^2 + ac^3 - ab^3 - ab^2c \]
We are going to subtract one from the other, so first cancel two terms duplicated between the two expressions:
\[ = 2ab^2c + b^2c^2 + bc^3 - a^3b - a^2bc \]
\[ = 2a^2bc + a^2c^2 + ac^3 - ab^3 - ab^2c \]

Place on one side by subtracting and set equal to zero:
\[ 2ab^2c + b^2c^2 + bc^3 - a^3b - a^2bc - 2a^2bc - a^2c^2 - ac^3 + ab^3 + ab^2c = 0 \]
Combine two pairs of like terms
\[ 3ab^2c + b^2c^2 + bc^3 - a^3b - 3a^2bc - a^2c^2 - ac^3 + ab^3 = 0 \]

Now the really tricky part.  Look for terms that are like $bx-ax$:
\[ bc^3 - ac^3 = (b-a)c^3 \]
\[ 3ab^2c - 3a^2bc = (b-a)3abc \]
\[ b^2c^2 - a^2c^2 = (b^2 - a^2)c^2 = (b-a)(b+a)c^2  \]
\[  ab^3 - a^3b = (b^2 - a^2)ab = (b-a)(b+a)ab \]

In other words the eight terms we had previously can be rewritten as
\[ (b-a) \ [ \ c^3 + 3abc + (b+a)c^2 + (b+a)ab \ ] \ = 0 \]
\[ (b-a) \ [ \ c^3 + 3abc + (b+a)(c^2 + ab) \ ] \ = 0 \]

And then, as we said above, we have the product of $(b-a)$ times something which is the sum of a bunch of positive terms.  The only way for that to be zero is if $b - a = 0$, and then $a = b$.

The two sides of $\triangle ABC$ are equal, it is an isosceles triangle.

$\square$

see \url{https://proofwiki.org/wiki/Steiner-Lehmus_Theorem}

\subsection*{afterward}

There is some interesting discussion in Coxeter as well.  According to what I can find on the web, most of the literature concerns the question of whether it is possible to provide a direct proof of the theorem.  The algebraic proof has been cited as such.

However, it depends on Stewart's Theorem, which as we derive it depends on the Law of Cosines, which depends in turn on the theorem of Pythagoras.  And although there are several hundred proofs of Pythagoras most (all?) of them depend on the sum of angles and also on the parallel postulate, which explicitly depends on a proof by contradiction.  

The question of a direct proof for Steiner-Lehmus is hard to answer conclusively.  I have a write-up from John Conway claiming that it is impossible, but I don't really understand his argument.  

\end{document}